\section{Modelo de calidad}
La calidad de la solución se asegura mediante controles técnicos, validación funcional y disciplina documental en todo el ciclo de vida.

\section{Estrategia de pruebas}
\begin{itemize}
  \item Pruebas unitarias de servicios críticos.
  \item Pruebas funcionales por flujo de negocio.
  \item Pruebas de regresion en funcionalidades sensibles.
  \item Pruebas de aceptacion con actores de negocio.
\end{itemize}

\section{Aseguramiento documental}
\begin{itemize}
  \item Control de versiones de especificaciones y manuales.
  \item Trazabilidad entre requerimientos, evidencia y resultados.
  \item Criterios de aprobación por etapa de liberación.
\end{itemize}

\section{Integracion continua y entrega continua (CI/CD)}
\begin{itemize}
  \item Gestión de entorno reproducible mediante control de dependencias y versionado de librerías.
  \item Automatización de validaciones técnicas en cada \textit{pull request} y en cada \textit{push} a rama principal.
  \item Ejecución automatizada de pruebas unitarias, chequeos de calidad de código y verificación de build.
  \item Publicación controlada con compuertas de aprobación para evitar regresiones en producción.
  \item Registro de evidencias de pipeline para auditoría de cambios y trazabilidad de entregas.
\end{itemize}

\section{Flujo de pipeline de referencia}
\begin{enumerate}
  \item Disparador de pipeline por evento de versionamiento (PR o push).
  \item Preparación del entorno y restauración de dependencias.
  \item Validación de estilo y calidad estática del código.
  \item Ejecución de suite de pruebas y consolidación de resultados.
  \item Generación de artefactos de build y reporte de evidencias.
  \item Promoción a entorno objetivo según política de aprobación.
\end{enumerate}

\section{Indicadores de calidad}
\begin{itemize}
  \item Tasa de incidencias por módulo.
  \item Cobertura de pruebas por componente.
  \item Tiempo medio de resolución de defectos.
  \item Cumplimiento de ventanas de liberación.
  \item Tasa de pipelines exitosos y estabilidad de despliegues.
\end{itemize}

\section{Mejora continua}
\begin{enumerate}
  \item Retroalimentacion de usuarios y areas responsables.
  \item Priorización de deuda técnica y funcional.
  \item Iteraciones de calidad con revisiones periódicas.
\end{enumerate}
